% STATUS: draft
% LAST_MODIFIED: 2026-07-26

\documentclass{internal-report}

\title{CUMCM 2021 C 题 — 初始分析报告}
\author{MathModel Team}
\date{2026年7月26日}

\begin{document}

\maketitle

\section{问题重述}

某建筑装饰板材生产企业每年按48周安排生产，需提前制定24周的原材料订购与转运计划。企业每周产能为2.82万立方米，消耗三种原材料：A类（0.6立方米/立方米产品）、B类（0.66立方米/立方米产品）、C类（0.72立方米/立方米产品）。A类和B类采购单价分别比C类高20\%和10\%，运输和储存单位费用三类相同。

企业从402家供应商采购原材料，原材料总体分A、B、C三类。供应商实际供货量可能多于或少于订货量，企业全部收购。需保持不少于满足两周生产需求的库存量。由8家转运商负责运输，每家运力为6000立方米/周，运输过程中存在损耗率。通常一家供应商每周的供货由一家转运商运输。

附件1提供402家供应商近5年（240周）的订货量与供货量数据。附件2提供8家转运商240周的运输损耗率数据。

\section{子问题拆解}

\subsection{子问题1：供应商重要性评估与筛选}
对402家供应商的供货特征进行量化分析，建立反映保障企业生产重要性的数学模型，筛选出50家最重要的供应商。

\textbf{关键挑战}：如何定义「保障企业生产重要性」？需综合考虑供货量、供货稳定性、品类分布等因素。

\subsection{子问题2：最少供应商与最优订购/转运方案}
确定满足生产需求的最少供应商数量，制定未来24周最经济的原材料订购方案，并制定损耗最少的转运方案。

\textbf{关键挑战}：订购与转运两层决策耦合；随机性（供货偏差+运输损耗）需要建模。

\subsection{子问题3：多目标成本压缩优化}
尽量多采购A类、尽量少采购C类原材料，同时希望转运损耗率尽量少。

\textbf{关键挑战}：三个目标存在冲突（A单价高但消耗率低），需要多目标优化方法建立Pareto前沿。

\subsection{子问题4：产能提升评估}
基于现有供应商和转运商条件，评估每周产能可提高多少，给出24周订购和转运方案。

\textbf{关键挑战}：需要在供应商供货能力和转运商运力双重约束下求解最大可行产能。

\section{相关知识}

\subsection{知识库已有覆盖}

\begin{itemize}
  \item \textbf{SP-006}：2021 CUMCM C题四种建模哲学对比 —— 直接对应本题，涵盖GA主导、双层随机规划、贪心启发式、LP序列四种方案。
  \item \textbf{MS-024}：熵权法+TOPSIS 多指标综合评价 —— 适用于子问题1的供应商评价。
  \item \textbf{MS-029}：投影寻踪法 —— 高维指标无监督降维评价，替代TOPSIS方案。
  \item \textbf{MS-030}：双层随机规划 —— 层级决策中分离目标优先级，适用于子问题2。
  \item \textbf{MS-031}：供货商时间模式分类 —— 订购建模前必须对供应商按时间序列特征分类。
  \item \textbf{MS-032}：多目标粒子群优化 —— 适用于子问题3的多目标权衡。
  \item \textbf{PF-008}：贪心算法供应商选择短视 —— 需注意避免局部最优。
  \item \textbf{PR-007}：订购建模前必须对供货商按时间模式分类。
  \item \textbf{AL-003}：遗传算法，AL-009：K-S检验，AL-017：动态权重贪心算法。
\end{itemize}

\subsection{参考论文（已解析）}
\begin{itemize}
  \item 2021C1：GA主导方案（6指标熵权+TOPSIS $\to$ GA统一求解订购+转运）
  \item 2021C2：双层随机规划（4指标熵权+TOPSIS $\to$ K-S检验供应商分类 $\to$ 双层随机规划）
  \item 2021C3：贪心启发式（9指标熵权+AHP修正 $\to$ 平稳/周期分类 $\to$ (s,S)存贮+双层贪心）
  \item 2021C4：LP序列（8指标投影寻踪+AGA $\to$ 周期/非周期异常值 $\to$ 线性规划+多目标PSO）
\end{itemize}

\subsection{知识空白}
\begin{itemize}
  \item \textbf{GAP-016}：最少供应商数量选择方法的系统性对比
  \item \textbf{GAP-017}：产能提升估计的验证方法论
\end{itemize}

\section{初始建模假设}

\begin{enumerate}
  \item 企业按恒定周产能（2.82万立方米）生产，不考虑产能波动。
  \item 供应商的历史供货模式在未来24周延续（周期性/统计分布可外推）。
  \item 转运商的历史损耗率模式在未来24周延续。
  \item 企业全部收购供应商实际供货量（超过订货量也不退还）。
  \item 原材料A、B、C之间可以按消耗率相互替代（通过调整产品结构）。
  \item 一家供应商每周只由一家转运商运输（软约束，尽量满足）。
\end{enumerate}

\section{待解决的开放问题}

\begin{enumerate}
  \item 子问题1的指标选择：除供货量均值、变异系数外，是否需考虑供货满足率、缺货频率、最大单周供货量等指标？
  \item 最少供应商数量：应以「稳定供货能力」还是「峰值供货能力」为判断标准？这直接影响结果（D1的39家 vs D3的16家 vs D4的24家）。
  \item 子问题2中订购与转运的耦合方式：是否可用双层规划统一求解？还是分步求解更稳健？
  \item 子问题3中三个目标的权重如何设定？各目标的量纲（成本、损耗率）不同，如何归一化？
  \item 子问题4中产能提升的约束瓶颈在哪里？是供应商供货极限还是转运商运力？
  \item 是否需要对随机变量（供货偏差、损耗率）进行分布拟合与K-S检验？
\end{enumerate}

\end{document}
